%% ============================================================================
%% Resumen (máx. 250 palabras según template oficial)
%% ============================================================================

\resumen{%
Este trabajo presenta el diseño y construcción de un robot colaborativo (cobot) con capacidad de agarre mediante ventosa, utilizando servomotores Herkulex DRS-0602 y el gripper piCOBOT. El proyecto se desarrolló en el marco del Tecnólogo en Mecatrónica de UTEC ITR Suroeste, con el objetivo de impulsar la adopción de robótica colaborativa y fortalecer el aprendizaje práctico. Se analizaron los aspectos mecánicos, eléctricos y electrónicos necesarios para el diseño; se desarrolló una librería propia para el control de los motores DRS-0602 ante la ausencia de soporte adecuado; y se construyó un prototipo de 4 grados de libertad mediante impresión 3D en PLA. La estructura modular, el sistema de comunicación UART y el diseño colaborativo (bordes redondeados, limitación de velocidades) permitieron validar el funcionamiento del cobot como plataforma demostrativa. Los resultados confirman la viabilidad de integrar componentes comerciales accesibles en un brazo robótico colaborativo de escritorio, con potencial para futuras mejoras en materiales, control embebido e interfaz de usuario.

\textbf{Palabras clave:} cobot; robótica colaborativa; Herkulex DRS-0602; piCOBOT; impresión 3D; industria 4.0.
}
